% Statement and cited-input packet. Target proof is separately embargoed.
% Frozen II source SHA256: 74497bfe327fae25c9cdcb17334b3bf42a9793a9df7c0a4404d97568f2f760f0
% Target owner: push:src:E.5.3; tag: 033P

\begin{theorem}[classification on the specified parameter domain]\hypertarget{tag-033P}{}\label{push:src:E.5.3}


There are $\epsilon>0$ and a real constant $C$ such that the
source construction, with its fixed logarithmic vector $\ell$,
defines a proper holomorphic family
$X_{\mu,c}$ over
$\{ |\mu|<\epsilon\}\times\{\operatorname{Im}c<C\}$.
It has the singular fibres of \ref{push:src:E.5.1} and the markings above.
Two of its members are biholomorphic if and only if
\[
 \mu'=\mu,\qquad c'-c\in2\mathbb Z.
\]
Every such biholomorphism covers the identity of the elliptic
base.  The quotient parameter in the second variable is
$v=e^{-\pi i c}$, with $\lambda=v^2$.
On a sufficiently small bidisc centred at any
$(0,\lambda_0)$ in Theorem 1.1 of \cite{ViaTwo26push}, choose
one logarithm of $\lambda$.  Within that actual bidisc,
\[
 X_{\mu,\lambda}\simeq X_{\mu',\lambda'}
 \quad\Longleftrightarrow\quad
 (\mu,\lambda)=(\mu',\lambda').
\]
\end{theorem}

% Cited input push:src:E.5.1; statement or labelled formula only.
\begin{proposition}[the input configuration]\hypertarget{tag-033N}{}\label{push:src:E.5.1}


The proposed configuration with three rank-one reduced singular fibres
does not describe the construction of \cite{ViaTwo26push}. The elliptic surface
$y^2=x^3+tx+1$ has fibre $III^*$ at infinity.  The constructed
threefold has two rank-one fibres at $p_1,p_2$, a rank-two
hexagonal fibre at $p_3$, and the multiple bielliptic fibre
$4B_*$ at infinity.
\end{proposition}

% Cited input push:src:E.5.2; statement or labelled formula only.
\begin{proposition}[monodromy, centralizer, and integral lattices]\hypertarget{tag-033O}{}\label{push:src:E.5.2}


Use the source's left-action marking
$H^1(F,\mathbb Z)=\mathbb Z\langle\psi,u,w,\delta\rangle$
and the ordered meridians $*,1,2,3$ with product one.
The cohomological matrices are
\[
\begin{aligned}
 T_*&=\begin{pmatrix}1&0&-2&-1\\0&0&-1&0\\0&1&0&-1\\0&0&0&1\end{pmatrix},&
 T_1&=\begin{pmatrix}1&0&2&2\\0&1&1&1\\0&0&1&0\\0&0&0&1\end{pmatrix},\\
 T_2&=\begin{pmatrix}1&0&0&0\\0&1&0&0\\0&-1&1&0\\0&0&0&1\end{pmatrix},&
 T_3&=\begin{pmatrix}1&0&0&-1\\0&1&1&0\\0&0&1&0\\0&0&0&1\end{pmatrix}.
\end{aligned}
\]
On ordinary dual homology put $H_i=(T_i^{-1})^t$,
$\ell=(1,-1,-1,0)^t$ and $\widehat\delta=(0,0,0,1)^t$.
Then
\[
 \operatorname{Cent}_{M_4(\mathbb Q)}(H_*,H_1,H_2,H_3)
 =\mathbb Q I\oplus\mathbb Q E_{41},
\]
and the cohomological centralizer is
$\mathbb Q I\oplus\mathbb Q E_{14}$.
The fixed and norm lattices at the multiple fibre are
\[
 \begin{aligned}
 \Lambda^{H_*}&=\mathbb Z\ell\oplus\mathbb Z\widehat\delta,\\
 N_4\Lambda&=4\mathbb Z\ell\oplus2\mathbb Z\widehat\delta,\\
 \Lambda^{H_*}/N_4\Lambda&\simeq\mathbb Z/4\oplus\mathbb Z/2.
 \end{aligned}
\]
The primitive vanishing lattices at the three reduced fibres are
\[
 K_1=\mathbb Z(\widehat w+\widehat\delta),\qquad
 K_2=\mathbb Z\widehat u,\qquad
 K_3=\mathbb Z\widehat w\oplus\mathbb Z\widehat\delta.
\]
Their invariant lattices are
\[
 \begin{split}
 \Lambda^{H_1}&=\mathbb Z(\widehat\psi-2\widehat u)
       \oplus\mathbb Z\widehat w\oplus\mathbb Z\widehat\delta,\\
 \Lambda^{H_2}&=\mathbb Z\widehat\psi\oplus
       \mathbb Z\widehat u\oplus\mathbb Z\widehat\delta,\\
 \Lambda^{H_3}&=K_3.
 \end{split}
\]
\end{proposition}

